\section{ウェーバーの法則に基づく感覚刺激}
ウェーバーの法則\cite{laming2009weber}によると，刺激強度の変化に気づくために必要な最小の変化量（弁別閾）は，元の刺激強度に比例する．

\begin{equation}
\frac{\Delta I}{I} = k
\end{equation}

ここで，$I$は基礎となる刺激の強さ，$\Delta I$は識別可能な最小の刺激変化量（丁度可知差異），$k$は定数（ウェーバー比）である．
一般に人間の感覚受容器における刺激受容応答は，物理的な刺激強度の対数に比例して知覚される特性（フェヒナーの法則）を持つが，これはウェーバーの法則を基礎としている．

Waldらは，ターゲットとの距離に応じてVRコントローラの振動強度を変化させることで，特別な訓練なしでも30度未満の精度で位置特定が可能であることを示した\cite{wald2025spatial}．彼らの研究では，振動強度の決定には，ウェーバーの法則に基づき，距離に対して指数関数的に減衰するモデルを採用している．

以上の知見を踏まえ，本研究ではウェーバーの法則に基づき，コントローラの位置と音源オブジェクトとの距離に応じて，振動強度を動的に調整するシステムを設計した．

\section{視聴覚刺激と音の高さ}
Millerらの研究では，視覚と聴覚の二感覚的な手がかりが利用可能な場合，対象の検出，識別，位置特定がより迅速かつ確実に行われることが示されている\cite{miller1982divided}．この知見に基づけば，音楽教育においても，聴覚情報と視覚情報を空間的に統合することで，学習効果を向上させられる可能性がある．

具体的には，視覚刺激と音程の関連性に関して，空間的に高い位置にある対象に対して高い音を，低い位置にある対象に対して低い音を対応させた場合に反応速度が向上する現象は，Spatial-Musical Association of Response Codes (SMARC) 効果として知られている\cite{rusconi2006spatial}．これは，音高が空間的な位置情報として表象され，垂直方向への空間マッピングが存在することを示唆している．

また，聴覚刺激に関して，PedersenらはVR空間におけるステレオ音響を利用した聴音トレーニングを実施し，空間認識が音楽経験者の音程認識に影響を与えることを明らかにした\cite{pedersen2020spatialized}．彼らの実験では，モノラル音源と比較してステレオ音源を使用した条件において，音楽経験者の音程識別テストの得点が有意に高かったことから，音楽経験者におけるステレオ音源を用いたトレーニングの有効性が示されている．

\section{VRと聴音トレーニング}
\ref{sec:background}で述べたように，音楽学習者が目指すべきはインナーヒアリングの獲得であり、音程認識能力の向上が最も基礎的な手段である\cite{woody2012playing}．この手段を満たすために，多くの研究で用いられているのがピッチインターバルである\cite{pedersen2020spatialized, fletcher2019virtual}。

本研究における聴音トレーニングでは，「ピッチインターバル（音程識別）」を採用する．McClaskeyらは，ピッチインターバルを用いたテストにより，音程識別能力には聴音トレーニングを受けた音楽経験者と受けていない音楽経験者では前者の方がスコアが高い傾向があった．この結果から，音程認識には聴音トレーニングのような特別な訓練を必要とする可能性を示唆しており，その習得には一定の訓練が必要であると述べている\cite{mcclaskey2017standard}．

また，音楽教育ソフトウェア「EarMaster」に搭載されている「Pitch Interval」は，2音間のピッチ距離を識別させるための基礎訓練項目である．音程は旋律（メロディ）を理解するための最小の構成単位であり，和音の響きを決定づける要素でもある．
